% FIG ch09-01 panel argument map (English)
\begin{tikzpicture}[
  x=1mm, y=1mm,
  font=\sffamily\small,
  text=BookInk,
  mainbox/.style={
    draw=BookTeal, line width=0.6pt, rounded corners=2pt,
    fill=BookPale, align=center, text width=150mm, inner sep=3mm,
  },
  explain/.style={
    draw=BookRule, fill=BookRule!35, rounded corners=2pt,
    align=left, text width=38mm, inner sep=2.5mm, minimum height=18mm,
  },
  moved/.style={
    draw=BookInk!40, dashed, rounded corners=2pt,
    align=left, text width=155mm, inner sep=2.5mm,
  },
]
  \node[mainbox] (q) at (87, 58)
    {Main question: does interface passivation raise both initial efficiency and operational stability?};
  \foreach \i/\x/\t in {
    a/20/schematic, b/62/morphology, c/104/performance, d/146/stability} {
    \node[explain] (\i) at (\x, 32) {
      \textbf{\i}\par
      Answers: what changed / layer continuity / initial efficiency / stability over time\par
      \scriptsize Evidence: \t
    };
  }
  \draw[-{Stealth[length=2mm]}, BookInk!60, thick] (q.south) -- (a.north);
  \draw[-{Stealth[length=2mm]}, BookInk!60] (a.east) -- (b.west);
  \draw[-{Stealth[length=2mm]}, BookInk!60] (b.east) -- (c.west);
  \draw[-{Stealth[length=2mm]}, BookInk!60] (c.east) -- (d.west);
  \node[moved] at (87, 8) {
    Candidate moved out (\textbf{example}): XRD pattern $\rightarrow$ Supplementary Information (does not serve the main question)
  };
\end{tikzpicture}
